\chapter{Cơ sở lý thuyết}
\label{ch:cosoly}

\section{Bài toán phát hiện đối tượng và chỉ số đánh giá}

Phát hiện đối tượng vừa định vị (bounding box) vừa phân loại vật thể. Chỉ số chuẩn là
\textbf{mean Average Precision} (mAP): trung bình của Average Precision trên các lớp. Hai biến
thể thường dùng là mAP@0,5 (ngưỡng IoU 0,5) và mAP@0,5:0,95 (trung bình trên 10 ngưỡng IoU từ
0,5 đến 0,95, nghiêm ngặt hơn về độ khớp hộp). Precision và Recall bổ sung góc nhìn về tỉ lệ dự
đoán đúng và tỉ lệ bỏ sót.

\section{YOLOv8 như một công cụ đo}

YOLOv8 \cite{yolov8} là mô hình phát hiện một giai đoạn, dự đoán trực tiếp hộp và lớp trên lưới
đặc trưng đa tỉ lệ. Trong đề tài này, một YOLOv8s đã huấn luyện được coi là \emph{công cụ đo cố
định}: ta không quan tâm cải tiến kiến trúc, mà dùng chính các lỗi và điểm số nó sinh ra làm dữ
liệu để phân tích độ tin cậy.

\section{Rò rỉ dữ liệu và trùng lặp}

Rò rỉ dữ liệu (\textit{data leakage}) xảy ra khi thông tin của tập kiểm thử ``lọt'' vào quá
trình huấn luyện. Với dữ liệu ảnh cắt từ video, một dạng data leakage phổ biến là \textbf{ảnh trùng
hoặc gần trùng xuất hiện ở nhiều split}. Perceptual hash (perceptual hash, ví dụ dHash) cho phép
phát hiện ảnh gần trùng: hai ảnh có mã băm giống hoặc rất gần nhau về khoảng cách Hamming được
coi là trùng tri giác. Nếu không loại trùng ở tập đánh giá, mô hình có thể ``ghi nhớ'' ảnh đã
thấy lúc huấn luyện, làm mAP cao hơn hiệu năng thực.

\section{Đo mất cân bằng lớp}

Ba thước đo bổ sung cho nhau được dùng: \emph{imbalance ratio} (tỉ số lớp nhiều nhất trên lớp ít
nhất), \emph{entropy chuẩn hoá} (bằng 1 khi phân bố đều) và \emph{hệ số Gini} về bất bình đẳng
(0 = cân bằng hoàn toàn; càng lớn càng tập trung). Cần phân biệt hệ số Gini bất bình đẳng với
Gini impurity trong cây quyết định --- đây là hai khái niệm khác nhau.

\section{Kiểm định thống kê sử dụng}

\textbf{Kiểm định Mann--Whitney U} \cite{mannwhitney} là kiểm định phi tham số so sánh phân phối
của hai nhóm độc lập mà không giả định phân phối chuẩn --- phù hợp khi so hiệu năng nhóm đèn (2
lớp) với nhóm biển (13 lớp). \textbf{Tương quan hạng Spearman} $\rho$ \cite{spearman} đo mức độ
đơn điệu giữa hai biến; ta dùng nó
để kiểm tra số lượng đối tượng của một lớp (độ hiếm) có liên hệ đơn điệu với hiệu năng lớp đó
hay không. Giá trị $\rho$ gần 0 với $p$ lớn nghĩa là không có tương quan.

\section{Định nghĩa kích thước vật thể}

Theo quy ước COCO \cite{coco}, vật thể được xếp \emph{small} nếu diện tích nhỏ hơn $32\times32$
px. Vật thể nhỏ khó phát hiện vì ít điểm ảnh mang thông tin; tỉ lệ vật thể nhỏ cao là một tín
hiệu khó khăn cho detector, đặc biệt ở độ phân giải đầu vào thấp.
